% !TeX root = ../sections/03_solutions.tex
\subsection{1.1.B.2.2}
\begin{proof}
    次に，
    \[
        \lim_{n\to\infty}
        \frac{n^2+2n}{n+1}
        =
        \infty
    \]
    を示す。

    正の無限大への発散を示すには，任意の $M>0$ に対して，
    ある自然数 $N$ が存在して，$n\geq N$ ならば
    \[
        \frac{n^2+2n}{n+1}>M
    \]
    となることを示せばよい。

    まず式を変形する。
    \[
        \frac{n^2+2n}{n+1}
        =
        \frac{(n+1)^2-1}{n+1}
    \]
    であるから，
    \[
        \frac{n^2+2n}{n+1}
        =
        n+1-\frac{1}{n+1}
    \]
    となる。

    ここで，$n\geq 1$ のとき
    \[
        0<\frac{1}{n+1}<1
    \]
    である。したがって
    \[
        n+1-\frac{1}{n+1}>n
    \]
    が成り立つ。

    つまり，
    \[
        \frac{n^2+2n}{n+1}>n
    \]
    である。

    任意に $M>0$ を取る。
    アルキメデスの性質より，
    \[
        N>M
    \]
    を満たす自然数 $N$ を取ることができる。

    このとき，$n\geq N$ ならば
    \[
        n\geq N>M
    \]
    である。

    したがって
    \[
        \frac{n^2+2n}{n+1}>n\geq N>M
    \]
    となる。

    よって，任意の $M>0$ に対して，ある自然数 $N$ が存在して，
    $n\geq N$ ならば
    \[
        \frac{n^2+2n}{n+1}>M
    \]
    が成り立つ。

    したがって
    \[
        \lim_{n\to\infty}
        \frac{n^2+2n}{n+1}
        =
        \infty
    \]
    である。
\end{proof}